\subsection{Related work}
Griggs and Yeh developed foundational results for $L(2,1)$-labeling
\cite{griggs}; Calamoneri surveys $L(h,k)$-labeling across graph classes
and parameter regimes \cite{calamoneri}. These results provide both
motivation and reference values for computational methods. Specialized
algorithms can exploit structure: $L(2,1)$-labeling on trees admits a
linear-time algorithm \cite{hasunuma}, whereas the tree complexity
result of Fiala, Golovach and Kratochv\'il includes the nondivisible
parameter regime represented here by $(3,2)$ \cite{fiala2008trees}.
Worst-case hardness does not imply that each sampled tree is difficult.
Murugan studied $L(2,1)$-labeling of corona products
\cite{murugan2015corona}; these families are established objects of study,
not new labeling classes introduced by this experiment. Work on structural
parameters motivates reporting graph structure alongside vertex count
\cite{fiala2018parameterized}; its parameterized algorithm is not evaluated here.

Exact SAT and ILP approaches precede this study. Shao and Vesel presented
assignment ILP and direct SAT formulations for distance-constrained
$L(k_1,\ldots,k_r)$-labeling, with experiments on $L(3,2,1)$
\cite{shao2013distance}. Their $L(2,1)$ work used direct SAT encoding on
strong products \cite{shao2014strong}. Tamura et al.\ describe the
order encoding of finite-domain comparisons using threshold variables
\cite{tamura2009order}. We use this representation for two configurable
distance constraints and evaluate the resulting implementation against
assignment ILP on common instances. The cited implementations have not
been reproduced under our protocol, so no performance ranking against
them is inferred.

Related variants provide context rather than interchangeable benchmark
values. Edge labeling corresponds to vertex labeling on a line graph,
subject to its distance convention \cite{berthe2023edge};
$L(h,1,1)$-labeling also constrains distance three \cite{king2010trees}.
Neither is the present two-distance model. The radio-labeling
preprint by Vu Thanh et al.\ uses order encoding and incremental SAT
\cite{vu2026radio}, while Sarhan et al.\ study a genetic algorithm for
radio labeling \cite{sarhan2026genetic}. Radio labeling with parameter
$r$ requires separation $r+1-d(u,v)$: $r=2$ gives $L(2,1)$, but $r=3$
also constrains distance three and therefore differs from $L(3,2)$.
Their numerical results are not reference values for our general
two-distance problem. An adapted genetic algorithm remains a possible
feasible-solution baseline and is not evaluated here.

Theoretical bounds, heuristic constructions and exact solvers provide
complementary evidence. A feasible labeling gives an upper bound; a
matching proven lower bound certifies optimality, regardless of how the
labeling was found. Otherwise, an exact solver must exclude smaller
spans to close the interval. This distinction underlies our separate
reporting of timed solver outcomes and combinatorial certificates.
